\documentclass[12pt,twoside]{book}

\usepackage[utf8]{inputenc}
\usepackage[T1]{fontenc}
\usepackage[polish]{babel}
\usepackage{lipsum}
\usepackage{graphicx}
\usepackage{float}
\usepackage{amsmath}
\usepackage{tikz}

\title{Skład tekstu w praktyce}
\author{Piotr Książek}
\date{\today}

\begin{document}

\maketitle
\tableofcontents
\listoffigures
\listoftables

% zad 3
\part{Podstawy LaTeX}

\chapter{Wprowadzenie}
\section{Pierwsza sekcja}
\lipsum[1]
\subsection{Pierwszy podrozdział}
\lipsum[2]

% zad 5-11
\section{Krój pisma}

\textbf{Tekst pogrubiony}

\textsl{Tekst pochylony}

\textit{Tekst kursywą}

\underline{Tekst z podkreśleniem}

\verb!Tekst maszynopis!

{\Large Większa czcionka}

{\small Mniejsza czcionka}

% zad 12
\section{Lista numerowana}
\begin{enumerate}
\item Pierwszy
\item Drugi
\item Trzeci
\end{enumerate}

% zad 13
\section{Lista wypunktowana}
\begin{itemize}
\item Jabłko
\item Gruszka
\item Banan
\end{itemize}

% zad 14
\chapter{Rysunki i tabele}
\section{Rysunek}
\begin{figure}[H]
\centering
\includegraphics[width=0.5\textwidth]{obrazek.png}
\caption{Przykładowy rysunek}
\label{fig:rysunek}
\end{figure}

% zad 15
\section{Tabela}
\begin{table}[H]
\centering
\caption{Tabela testowa}
\label{tab:test}
\begin{tabular}{|c|c|c|}
\hline
1 & 2 & 3 \\ \hline
4 & 5 & 6 \\ \hline
7 & 8 & 9 \\ \hline
10 & 11 & 12 \\ \hline
13 & 14 & 15 \\ \hline
\end{tabular}
\end{table}

% zad 16 i 17
\chapter{Wzory}
\section{Pole koła}
\begin{equation}
P = \pi r^2
\label{eq:kolo}
\end{equation}

\section{Równanie kwadratowe}
\begin{equation}
\Delta = b^2 - 4 \cdot a \cdot c
\label{eq:delta}
\end{equation}
\begin{equation}
x_1 = \frac{-b - \sqrt{\Delta}}{2 \cdot a}
\end{equation}
\begin{equation}
x_2 = \frac{-b + \sqrt{\Delta}}{2 \cdot a}
\end{equation}

% zad 18
\section{Odwołania}
Rysunek \ref{fig:rysunek} jest na stronie \pageref{fig:rysunek}
Tabela \ref{tab:test} ma 3 kolumny
Wzór \ref{eq:kolo} to pole koła a wzór \ref{eq:delta} to delta

% zad 19
\chapter{Przypisy i znaki}
Zdanie z przypisem\footnote{Przypis dolny}

% zad 20
\section{Znaki specjalne}
\$ \# \% \^{} \& \textquotedbl{} [ ] \textbackslash{} \textasciitilde{} \{ \} | \_ @

% zad 22
\chapter{Grafika TikZ}
\section{Rysunki w PGF/TikZ}

\begin{figure}[H]
\centering
\begin{tikzpicture}[node distance=1.25cm, main/.style={draw, circle}]
\node[main] (1) {a};
\node[main] (2) [below right of=1] {b};
\node[main] (3) [above right of=1] {c};
\node[main, draw=none] (4) [right of=3] {~};
\draw[->] (1) to (2);
\draw[->, color=green] (1) to (3);
\draw[->, dashed, color=blue] (3) to (4);
\end{tikzpicture}
\caption{Początkowy graf}
\end{figure}

\begin{figure}[H]
\centering
\begin{tikzpicture}[node distance=1.25cm, main/.style={draw, circle}]
\node[main] (1) {a};
\node[main] (2) [below right of=1] {b};
\node[main] (3) [above right of=1] {c};
\node[main, draw=none] (4) [right of=3] {~};
\node[main] (5) at (4.2,0) {$\Delta$};
\node[main] (6) at (6.2,-0.88) {z};
\draw[->] (1) to (2);
\draw[->, color=green] (1) to (3);
\draw[->, dashed, color=blue] (3) to (4);
\draw[dotted] (1) to (5);
\draw[->] (3) to (5);
\draw[->] (2) to (5);
\draw[red] (2) to (6);
\draw[->>, red] (6) to (5);
\end{tikzpicture}
\caption{Końcowy graf}
\end{figure}

\end{document}